% Ref：https://arxiv.org/pdf/2505.02012

% 1. LLM-based testing.
% \subsection{LLM-based Testing}
% % {\bf \underline{LLM-based Testing.}}
% The advent of Large Language Models (LLMs) has introduced new paradigms for automated software testing. A significant body of work explores the use of LLMs as direct, online test case generators. For instance, Fuzz4All~\cite{fuzz4all} is a universal fuzzing framework that leverages LLM-driven prompt generation and mutation to create diverse test inputs for a wide variety of programming languages. In more specialized domains, WhiteFox~\cite{YangDLY0J024} applies LLMs to white-box compiler fuzzing. It employs a sophisticated two-model approach where one LLM analyzes compiler source code to understand optimization passes, and a second LLM generates targeted test programs designed to exercise specific compiler behaviors. Similarly, KernelGPT~\cite{YangZZ25} enhances OS kernel fuzzing by prompting LLMs to synthesize system call specifications from natural language documentation, which are then used to guide traditional fuzzing engines. While powerful, these direct-generation methods often face challenges related to low throughput and high computational costs, which can limit their applicability in high-frequency or resource-constrained testing environments.

% An alternative strategy involves using LLMs in an offline capacity to augment or create assets for traditional testing tools. This approach separates the costly LLM inference phase from the high-throughput execution phase. MetaMut~\cite{OuL0024} exemplifies this paradigm by using an LLM to automatically generate high-quality, domain-specific mutators for compiler fuzzing. These generated mutators are then integrated into a conventional fuzzing framework, thereby combining the creative, knowledge-driven capabilities of LLMs with the speed and efficiency of traditional methods. This body of work demonstrates that LLMs are being effectively explored both as direct test generators and as powerful offline assistants that enhance existing testing infrastructure, each with distinct trade-offs in performance, cost, and flexibility.


% \subsection{Test Case Generation for DBMSs}

% % {\bf \underline{Test Case Generation for DBMSs.}}
% Automated test case generation is critical for ensuring the correctness and robustness of complex Database Management Systems (DBMSs). The primary methods in this domain can be broadly classified as mutation-based or generation-based. Mutation-based testing involves modifying a corpus of existing SQL statements, or seeds, to produce new and varied test cases. Griffin~\cite{FuLWWJ22} advances this technique with a grammar-free approach, using a lightweight metadata graph of the database state to guide mutations and improve the semantic validity of the resulting queries. The scope of mutation has also expanded beyond traditional SQL; BuzzBee~\cite{YangCZCL24} extends fuzzing to cover multiple database models, while Sedar~\cite{Fu0W024} improves input diversity by developing techniques to transfer and adapt SQL seeds across different DBMS dialects. These approaches are highly effective at discovering crash, hang, and assertion failure bugs due to their ability to generate a high volume of syntactically diverse inputs.

% In contrast, generation-based testing constructs queries from scratch using predefined rules and grammars. SQLsmith~\cite{SQLsmith} is a foundational tool in this area, capable of generating random yet syntactically well-formed SQL queries by leveraging database schema information. The SQLancer framework~\cite{SQLancer1,SQLancer2,SQLancer3} represents a significant advancement in this category, employing hand-crafted, dialect-specific generators in conjunction with sophisticated test oracles (e.g., Ternary Logic Partitioning and Non-optimizing Reference Engine) to effectively detect subtle logic bugs. The primary advantage of generation-based methods is their fine-grained control over query structure and semantics, which is crucial for satisfying the strict constraints of logic bug oracles. The field thus presents a clear trade-off between mutation-based methods, valued for their speed and effectiveness in finding memory-related errors, and generation-based methods, prized for their semantic precision in uncovering complex logical inconsistencies.

% % 2. Test case generation for DBMSs.
\label{sec:related}
\subsection{LLM-based Testing}

The advent of LLMs has introduced new paradigms for automated software testing. Existing LLM-based testing approaches can be divided into two categories: direct generation and offline augmentation. For direct generation, Fuzz4All~\cite{fuzz4all} leverages LLM-driven prompt generation and mutation to create diverse test inputs for multiple programming languages. WhiteFox~\cite{YangDLY0J024} applies a two-model approach for white-box compiler fuzzing, where one LLM analyzes compiler source code and another generates targeted test programs. KernelGPT~\cite{YangZZ25} enhances OS kernel fuzzing by prompting LLMs to synthesize system call specifications from documentation. For offline augmentation, MetaMut~\cite{OuL0024} uses LLMs to automatically generate domain-specific mutators for compiler fuzzing, separating the costly LLM inference from high-throughput execution. Different from these approaches that focus on general-purpose programming languages, compilers, or operating systems, our work targets DBMS testing, where the challenge lies in adapting to proprietary SQL dialects and achieving high code coverage with lightweight local models under industrial resource constraints.

\subsection{Test Case Generation for DBMSs}

Existing methods for generating test cases for DBMSs can be divided into two categories: mutation-based and generation-based testing. Mutation-based approaches modify existing SQL statements to produce new test cases. Sedar~\cite{Fu0W024} improves input diversity by transferring SQL seeds across DBMS dialects. Griffin~\cite{FuLWWJ22} uses a grammar-free approach with a lightweight metadata graph to guide semantically-correct query mutations. Generation-based approaches construct queries from scratch using predefined rules. SQLsmith~\cite{SQLsmith} generates random yet well-formed SQL queries by leveraging schema metadata. SQLancer~\cite{SQLancer1,SQLancer2,SQLancer3} employs hand-crafted dialect-specific generators with sophisticated test oracles to detect logic bugs. Different from these traditional approaches that rely on manual grammar engineering or random mutation, our work is the first to leverage lightweight LLMs with hierarchical feature guidance and MCTS-based mutation guided by coverage feedback to systematically improve code coverage.